%% LyX 2.4.4 created this file.  For more info, see https://www.lyx.org/.
%% Do not edit unless you really know what you are doing.
\documentclass[english]{article}
\usepackage[T1]{fontenc}
\usepackage[latin9]{inputenc}
\usepackage{units}
\usepackage{enumitem}
\usepackage{amsmath}
\usepackage{amsthm}
\usepackage{amssymb}
\usepackage{geometry}
\geometry{verbose,tmargin=1in,bmargin=1in,lmargin=1in,rmargin=1in}
\PassOptionsToPackage{normalem}{ulem}
\usepackage{ulem}

\makeatletter
%%%%%%%%%%%%%%%%%%%%%%%%%%%%%% Textclass specific LaTeX commands.
\numberwithin{equation}{section}
\numberwithin{figure}{section}
\newlength{\lyxlabelwidth}      % auxiliary length 
\theoremstyle{plain}
\newtheorem{thm}{\protect\theoremname}[section]
\theoremstyle{definition}
\newtheorem{xca}[thm]{\protect\exercisename}

%%%%%%%%%%%%%%%%%%%%%%%%%%%%%% User specified LaTeX commands.
\usepackage{fancyhdr}
\usepackage{lastpage}
\pagestyle{fancy}
\usepackage{enumitem}
\usepackage{ifthen}
\everymath=\expandafter{\the\everymath\displaystyle}
%\DeclareMathSizes{10}{10}{10}{10}
\usepackage{multicol}

\usepackage{tikz}
\usetikzlibrary{patterns}
\usetikzlibrary{plotmarks}
\usepackage{pgfplots}
\pgfplotsset{compat=newest}

\usepackage{calc}
\usepackage[nomessages]{fp}

\usepackage{datetime}
% Added by lyx2lyx
\setlength{\parskip}{\smallskipamount}
\setlength{\parindent}{0pt}

\makeatother

\usepackage{babel}
\providecommand{\exercisename}{Exercise}
\providecommand{\theoremname}{Theorem}

\begin{document}
\renewcommand{\author}{Dr.\,James Ashe}

\newcommand{\classNum}{439}

\newcommand{\classSec}{A}

\newcommand{\nameType}{solutions}

\newcommand{\examType}{Homework 5}

\newcommand{\Date}{\formatdate{6}{2}{2014}} %%\AdvanceDate[12] \today

%\newdate{my_date}{\day}{\month}{\year}%   
%\AdvanceDate[-5] 
%\displaydate{my_date}

%%First page's header%%%%%%%%%%%%%%%%%%%%%%
\fancypagestyle{firststyle} %%header settings for first pagr
{
	\fancyhead[L]{MTH \classNum{}(\classSec{}) \author{}\\
	\textbf{\examType{}} \Date{}}
	\fancyhead[C]{\textbf{\nameType}}
	\setlength{\headsep}{14pt}
}
%%%%%%%%%%%%%%%%%%%%%%%%%%%%%%%%%%%%%%%%%%

%%Next page's header%%%%%%%%%%%%%%%%%%%%%%%
%\setlist{nolistsep} %eliminates space between list items
\lhead{MTH \classNum{}(\classSec{})} 
\chead{\textbf{\nameType}} 
\cfoot{\thepage{}~of~\pageref{LastPage}} 
\setlength{\headsep}{5pt} 
%%%%%%%%%%%%%%%%%%%%%%%%%%%%%%%%%%%%%%%%%%%

%%Various settings; see the preamble for other settings%%%%%
%%\setenumerate[0]{leftmargin=12pt,itemindent=0pt,labelwidth=0pt}
\renewcommand{\labelenumi}{\textbf{\arabic{enumi}.}}
\renewcommand{\labelenumii}{\textbf{\alph{enumii}.}}
\thispagestyle{firststyle}
%%%%%%%%%%%%%%%%%%%%%%%%%%%%%%%%%%%%%%%%%%

\null

\vspace{1cm}

\setcounter{section}{3}

\Large{\textbf{Chapter 3}} \large

\textbf{\textbf{\examType} due: \formatdate{10}{4}{2014}}

\setcounter{thm}{13}

All problems \uline{without }a $\bigstar$ from section 4.1; exercises
3.14-3.19. 
\begin{xca}
$\bigstar$
\end{xca}

\begin{xca}
Give two examples of elements in $\mbox{\ensuremath{\mathbb{Z}}}/4\mathbb{Z}=\mathbb{Z}/\left\langle 4\right\rangle $,
and add and multiply them.

\begin{proof}
Possible elements are $0+\left\langle 4\right\rangle =\left\langle 4\right\rangle $,
$1+\left\langle 4\right\rangle $, $2+\left\langle 4\right\rangle $,
$3+\left\langle 4\right\rangle $. 

\begin{itemize}
\item $a+\left\langle 4\right\rangle =\left\{ \dots,a-4,0+a,4+a,8+a,\dots\right\} $ 
\end{itemize}
\end{proof}
Take $1+\left\langle 4\right\rangle $ and $2+\left\langle 4\right\rangle $.
\[
\left(1+\left\langle 4\right\rangle \right)\oplus\left(2+\left\langle 4\right\rangle \right)=\left(1+2\right)+\left\langle 4\right\rangle =3+\left\langle 4\right\rangle =1+\left\langle 4\right\rangle 
\]
\[
\left(1+\left\langle 4\right\rangle \right)\otimes\left(2+\left\langle 4\right\rangle \right)=2+\left\langle 4\right\rangle 
\]

\end{xca}

\begin{xca}
Find all the distinct cosets of $\mbox{\ensuremath{\mathbb{Z}}}/4\mathbb{Z}$.

\begin{proof}
Recall that $\mathbb{Z}/\left\langle 4\right\rangle \cong\mathbb{Z}_{4}$
so it has $4$ elements. From above they are $0+\left\langle 4\right\rangle =\left\langle 4\right\rangle $,
$1+\left\langle 4\right\rangle $, $2+\left\langle 4\right\rangle $,
$3+\left\langle 4\right\rangle $. 

To see that they are distinct note that $4\nmid a-b$ for any $0\leq a,b\leq3$
since then $a-b<4$. 
\end{proof}
\end{xca}

\begin{xca}
Let $S=\left\{ a+bi\colon b\in2\mathbb{Z}\right\} $ ($b$ is even
and $i=\sqrt{-i}$). Prove that $S$ is a subring of $\mathbb{Z}[i]$
but that $S$ is not an ideal of $\mathbb{Z}[i]$. 

\begin{proof}
Let $a+2ni,a'+2mi\in S$ for some $a,a',n,m\in\mathbb{Z}$. Using
the subring test, 
\[
(a+2ni)-(a'+2mi)=(a-a')+2(n-m)i\in S
\]
 and 
\[
(a+2ni)(a'+2mi)=aa'+2ami+2ani-2nm=(aa'-2nm)+2(m+n)i\in S
\]
so $S$ is a ring.

To show that $S$ is not an ideal in $\mathbb{Z}[x]$ consider $1+2i\in S$
and $1+i\in\mathbb{Z}[x]$. 
\[
\left(1+2i\right)\left(1+i\right)=2+2i+i-2=0+3i\notin S\mbox{.}
\]
\end{proof}
\end{xca}

-

\begin{xca}
Let $S=\left\{ a+bi\colon b\in2\mathbb{Z}\right\} $ ($b$ is even
and $i=\sqrt{-i}$). Prove that $\mathbb{Z}[i]/S$ is not a ring.

\begin{proof}
By exercise 3.4 $S$ is not an ideal. So by theorem 4.3 $\mathbb{Z}[i]/S$
is not a ring. 
\end{proof}
\end{xca}

\begin{xca}
Prove that $\mathbb{R}[x]/\left\langle x^{2}+1\right\rangle $ is
a ring. (We know that the ring of polynomials with integer coefficients,
$\mathbb{R}[x]$, is a ring).

\begin{proof}
By theorem 3.4 we need only show that $\left\langle x^{2}+1\right\rangle $
is an ideal. If $g_{1}(x),g_{2}(x)\in\left\langle x^{2}+1\right\rangle $
then $g_{1}(x)=(x^{2}+1)f_{1}(x)$ and $g_{2}(x)=(x^{2}+1)f_{2}(x)$
for some $f_{1}(x),f_{2}(x)\in\mathbb{R}[x]$. Using the ideal test,
\begin{eqnarray*}
g_{1}(x)-g_{2}(x) & = & (x^{2}+1)f_{1}(x)-(x^{2}+1)f_{2}(x)\\
 & = & (x^{2}+1)\left(f_{1}(x)-f_{2}(x)\right)\in\left\langle x^{2}-1\right\rangle 
\end{eqnarray*}
and for any $h(x)\in\mathbb{R}[x]$,
\[
h(x)g_{1}(x)=g_{1}(x)h(x)=h(x)\left((x^{2}+1\right)f_{1}(x)=h(x)f_{1}(x)(x^{2}+1)\in\left\langle x^{2}+1\right\rangle 
\]
so $\left\langle x^{2}+1\right\rangle $ is an ideal which implies
that $\mathbb{R}[x]/\left\langle x^{2}+1\right\rangle $ is a ring
by the theorem. 
\end{proof}
\end{xca}

\begin{xca}
Give two examples of elements in $\mbox{\ensuremath{\mathbb{Z}}}[i]/\left\langle 2-i\right\rangle $,
and add and multiply them.

\begin{proof}
$1+\left\langle 2-i\right\rangle $ and $(2+i)+\left\langle 2-i\right\rangle $
are two elements in the factor ring $\mathbb{Z}[i]/\left\langle 2-i\right\rangle $.
\[
\left[1+\left\langle 2-i\right\rangle \right]\oplus\left[(2+i)+\left\langle 2-i\right\rangle \right]=(3+i)+\left\langle 2-i\right\rangle 
\]
and 
\[
\left[1+\left\langle 2-i\right\rangle \right]\left[(2+i)+\left\langle 2-i\right\rangle \right]=(2+i)+\left\langle 2-i\right\rangle .
\]
\end{proof}
\begin{itemize}
\item Note that $2-i\equiv0\bmod\left\langle 2-i\right\rangle $ so then
$2\equiv i$, and $(2+i)+\left\langle 2-i\right\rangle \equiv4+\left\langle 2-i\right\rangle $.
We can easily check this: $(2+i)-4=-2+i=-1(2-i)\in\left\langle 2-i\right\rangle $. 
\end{itemize}
\end{xca}

\begin{xca}
Find all the distinct cosets of $\mbox{\ensuremath{\mathbb{Z}}}[i]/\left\langle 2-i\right\rangle $.

\begin{proof}
Let $\left\langle 2-i\right\rangle \in\mathbb{Z}[i]$. $\mathbb{Z}[i]=\left\{ x+yi\colon x,y\in\mathbb{Z}\right\} $,
called the \emph{Gaussian integers. }For any two Gaussian integers
$z$ and $z'$ we have \emph{. $(2-1i)z-(2-i)z'=(2-i)(z-z')\in\left\langle 4-i\right\rangle $},
and $z_{1}\left(z_{2}(2-i)\right)=\left(z_{2}\left(2-i\right)\right)z_{1}=z_{1}z_{2}\left(2-i\right)\in\left\langle 2-i\right\rangle $.
So $\left\langle 2-i\right\rangle $ is an ideal of $\mathbb{Z}[i]$. 

By the theorem, $\mathbb{Z}[i]/\left\langle 4-i\right\rangle $ is
a ring with elements (cosets) of the form $x+yi+\left\langle 2-i\right\rangle $
for $x,y\in\mathbb{Z}$. When dealing with cosets in $\mathbb{Z}[x]/\left\langle 2-i\right\rangle $
we can treat $2-i$ as $0$ since $2-i+\left\langle 2-i\right\rangle =0+\left\langle 2-i\right\rangle $.
This allows us to simplify the set of \emph{distinct }cosets. For
example, since $2-i\equiv0$ we also know that $2\equiv i$ since
$2-i+i=2$. So then 
\[
6+3i+\left\langle 2-i\right\rangle =6+3(2)+\left\langle 2-i\right\rangle =12+\left\langle 2-i\right\rangle 
\]

$2\equiv i\Rightarrow4\equiv-1$ and $4+1=5\equiv0$. So we should
expect there to be $5$ distinct cosets, and

\[
\mathbb{Z}[i]/\left\langle 4-i\right\rangle =\left\{ 0+\left\langle 4-i\right\rangle ,1+\left\langle 4-i\right\rangle ,\dots,4+\left\langle 4-i\right\rangle \right\} 
\]
To see that this is the case, note that any other coset will be some
integer multiple of $1+\left\langle 2-i\right\rangle $. Since $5\left(1+\left\langle 2-i\right\rangle \right)=0+\left\langle 2-i\right\rangle $,
$5\left(1+\left\langle 2-i\right\rangle \right)$ has order $1$ or
$5$. If it is $1$, then $1+\left\langle 2-i\right\rangle =0+\left\langle 2-i\right\rangle $
and $1\in\left\langle 2-i\right\rangle $ so $1=\left(x+yi\right)\left(2-i\right)=2x-ix+2yi+y=2x+y+i(-x+2y)$
for some integers $x$ and $y$. Since $1\in\mathbb{R}$ this implies
that $1=2x+y$ and $0=-x+2y$. Solving this system gives $y=\nicefrac{1}{5}\notin\mathbb{Z}$,
a contradiction. So the order of $\mathbb{Z}[x]/\left\langle 2-i\right\rangle $
is $5$. From the operations we can also see that this ring is essentially
the same as $\mathbb{Z}_{5}$, a field. 
\end{proof}
\end{xca}


\end{document}
